% ARCHIVED at v3.90 (2026-09-24) from chapters/06_results.tex -- author: "keep refine each sub conclusion and the
% final 6.4", under Working_Space/GUIDE_20260924_results_storyline_author.md. Verbatim, not part of the build.

% ---- 6.1.4 D3IL-avoiding, first two paragraphs ----
The supported combinations are collected in \autoref{tab:avoiding-conclusion}. At the protocol of
\ac{DPCC}, CI-MeanFM at one network evaluation with temporal consistency and per-step projection
reaches $1.000$ success with constraint satisfaction in $59.2$ control steps and $18.1$\,ms per step,
and MeanFM at the same budget and projector $0.967$ in $58.6$ steps and $18.3$\,ms. The baseline of
record reaches $1.000$ in $70.1$ steps and $553.4$\,ms. Both flow-based rows retain the temporal U-Net
of $4.0$\,M parameters used by the diffusion baseline.

On D3IL-avoiding, CI-MeanFM and FM at one network evaluation dominate the baseline of record ---
constraint satisfaction held, fewer control steps, a thirtieth of the time per action --- and MeanFM at
the same budget trades one episode in thirty for the same saving. Among the flow-based models the step
budget and the selection rule set the ordering, not the flow objective.

% ---- 6.2.4 D3IL-aligning, two paragraphs ----
D3IL-aligning selects analytic average-velocity matching as the generative model: at twenty evaluations
it ends closer to the target than instantaneous-velocity matching and the diffusion baseline, whereas
consistency-interpolated average-velocity matching does not improve with the larger budget. It selects
endpoint projection as the projection method: at the large budget this task needs, endpoint projection
with random selection keeps all ten contexts free of violations and still moves the box in eight, at the
time per control step of per-step projection at $\nfe=20$ and at $37\,\%$ of it at $\nfe=10$. The
conclusion is limited to training-set contexts.

Without projection, analytic average-velocity matching closes $84\,\%$ of the distance at twenty
evaluations, against $14$, $10$ and $-2\,\%$ for the consistency-interpolated model, instantaneous-velocity
matching and the diffusion baseline, and the consistency-interpolated model reaches $60\,\%$ only at a
hundred. Under per-step projection with random selection, analytic average-velocity matching
Pareto-dominates the baseline's best projected row, and endpoint projection keeps more contexts free of
violations than per-step projection in seven of nine comparisons at the budget the task needs. The executed end-effector paths of
the three projection settings, context by context, are drawn in \autoref{app:aligning-paths}.

% ---- 6.3.1.3 UAV-pillars conclusion ----
\subsubsection{Conclusion}\label{sec:res:uav:pillars:conclusion}

UAV-pillars works by changing only the plant, from the manipulator to the quadrotor with its controller,
as far as the declared constraints go. Before projection the flown plans violate them as the executed ones
do, and after projection neither violates them at all, at the same control steps to within three and the
same time per action to within a millisecond. The reason is how the scene is built
(\autoref{sec:method:env:uav}): the planner, its plan, its observation, its projection and its constraint
rows are those of D3IL-avoiding, in that benchmark's coordinates, and only the plant that executes each
commanded position differs. The arena is the table under the similarity map
$X = s\,(y\sidx{a}-y_{0})$, $Y = -s\,(x\sidx{a}-x_{0})$ at a fixed altitude, with $s = 36$, and the
vehicle follows each mapped setpoint closely: its flown position stays within about $0.007$ of the
commanded one in the table's units, closer than the manipulator's own $0.02$ to $0.04$ on D3IL-avoiding
and well inside the tightening margin of $0.025$ that the projection keeps from every declared boundary.
The flown path therefore keeps the constraint satisfaction the projection gives the plan.
The result still differs, $0.600$ to $0.700$ success within the constraints against $0.967$ to $1.000$,
although the demonstrations, the planner and the projection are those of D3IL-avoiding and only the plant is
replaced; the difference is the collisions, and both reasons for it are the quadrotor's own. The first is
why they happen: the projection keeps the plan out of the declared set only, the plans pass the five pillars
the set leaves out as closely as the demonstrations passed the posts, at the rod's radius, and the scale
that makes that radius the rotor reach leaves the vehicle no margin there, so that where a plan skims a
pillar the rotors meet it --- on $50$ of the $51$ collisions the commanded path itself is inside the reach.
The second is what a collision costs: for a quadrotor it is catastrophic, a rotor strike that in most cases
ends in lost control, and the flight ends at the first contact. The declared constraint set transfers
exactly; freedom from collision with the obstacles the set leaves out does not.

Where UAV-pillars can test it, the relationship D3IL-avoiding established between the models and the
projection holds (\autoref{sec:res:avoiding:conclusion}). The two average-velocity models keep their order
and their gap: CI-MeanFM succeeds within the constraints on one flight in thirty more than analytic
average-velocity matching at both budgets, on UAV-pillars as on D3IL-avoiding, within the spread over the
seeds. Per-step projection is still what turns plans that violate a constraint in almost every episode into
plans that violate none, and the collisions do not re-order the two models: both lose the same
share, on the same geometry. What UAV-pillars does not test is the comparison with the baseline. Neither
the diffusion model nor instantaneous-velocity matching was flown, so the dominance of CI-MeanFM and FM
over the baseline of record remains a result of D3IL-avoiding.
\dataref{Setpoint gap: \texttt{DA\_20260923\_pillars\_v2\_live.md} \S2 --- the 95th percentile per flight of
the distance between the commanded and the flown position, mapped back to the table, $0.0066$ on average,
against $0.02$--$0.04$ for the manipulator's own gap on D3IL-avoiding; mean tracking error $0.18$\,m in the
arena. Tightening $\delta = 0.025$ (\autoref{tab:eval}). The map is \texttt{eq:method:env:pillarsmap} of
Ch 4 (\texttt{uav\_avoiding\_bridge/frame.py}); it is written out here rather than referenced because the
review bundle collapses Ch 4 and a reference into it prints \enquote{??}. The one-in-thirty gap:
$1.000$ against $0.967$ on D3IL-avoiding; $0.633$ against $0.600$ and $0.700$ against $0.667$ on UAV-pillars.}

% ---- 6.3.2.4 UAV-corridor conclusion ----
\subsubsection{Conclusion}\label{sec:res:uav:corridor:conclusion}

On UAV-corridor the generative model does not matter; the budget does. The demonstrations carry little
information --- straight lines on three lanes at one altitude per episode, a single path forward from any
position (\autoref{tab:uav-demos}) --- and the three flow-based objectives learn the same plans from them.
Without projection their flights cannot be told apart; with it they leave the same violating steps at a
shared budget and projection method, to within $1.6$ per flight, while under the tilt the budget takes them
from about ten at one evaluation to one or two at twenty. Against the baseline they reach the end of the
corridor on every flight at a twenty-fourth (tilt) and a ninth (hump) of its time per action at their
cheapest, and none leaves as few violating steps: at best $1.4$ and $0.2$ per flight, against its $0.8$ and
$0.0$.

Endpoint projection beats the per-step projection of \ac{DPCC} in one place: over the hump at twenty
evaluations with a single candidate, where it leaves fewer violating steps at less time per action, $0.2$
against $1.8$ at $291$ against $355$\,ms, on flights about thirty control steps longer. Under the tilt it
does not: its fewest violating steps, $1.4$ against $2.1$, cost $417$ against $360$\,ms, a trade-off. At
three and five evaluations it is as cheap or cheaper but leaves as many violating steps or more, and at one
and two it has no guiding step. Neither projector removes every violation: a residue stays where each
constraint ends.
\dataref{v3.88 (author: \enquote{too long, is the model no diff? only K matters? should yes, since the expert
data contain low info. how is the projection, did the HF beat the PCC? that is the RQ, plz answer clearly});
v3.87 (author: the key conclusion, the model does not matter on straight-line demonstrations). The tie:
\autoref{tab:uav-corridor-raw} and \autoref{tab:uav-corridor}, MeanFM / CI-MeanFM / FM at shared budgets, the
largest gap tilt per-step $\nfe=3$ $6.6/5.0/5.8$ (sd $0.8$); the full list in the v3.87 changelog. RQ2
(\autoref{sec:res:uav:projection}, FM): hump $\nfe=20$ endpoint, one candidate, $0.2\pm0.4$ violating steps,
$316.2$ control steps, $291.0$\,ms, against per-step $1.8\pm0.4$, $285.1$, $355.2$\,ms; tilt endpoint $1.4\pm0.7$
at $416.9$\,ms against per-step $2.1\pm0.7$ at $360.4$\,ms (one candidate: $2.3$ at $289$\,ms). Price:
$666.3/27.4 = 24$, $654.2/71.9 = 9.1$. The residue: \autoref{sec:res:uav:projected} (paths). The v3.86--v3.87
long form (altitude, the hand-over, the relation to D3IL-avoiding) is in \texttt{withheld/20260924\_v3.88\_archive/}.}

% ---- 6.3.3.4 UAV-s-curve conclusion ----
\subsubsection{Conclusion}\label{sec:res:uav:scurve:conclusion}

\provisional{Prefilled on the unprojected grid of \autoref{tab:uav-scurve} alone. The projected cells (R44b)
and the controller comparison (R44c) are pending, and this conclusion is rewritten when they land.}

\provisional{UAV-s-curve is the second scene whose demonstrations are simple. They follow one route,
straight segments joined by arcs of at most $0.45$\,m radius, shifted sideways by at most $0.04$\,m per
episode, and the route clears the constraint set by $0.121$\,m (\autoref{tab:uav-demos},
\autoref{fig:expert-uav}); as on the straight lanes of UAV-corridor, from any observed position the
demonstrations continue along a single path. On such data the average-velocity objectives are not ahead at
the budget they are built for. At one evaluation instantaneous-velocity matching crosses the finish line on
nine flights of ten, the consistency-interpolated model on six and analytic average-velocity matching on
three, and for every flow-based model the crossings fall and the inversions rise as the budget grows from
one evaluation to twenty; the diffusion baseline crosses on none. No flight is free of violations: the
flow-based flights follow their plans with a mean tracking error of $0.30$ to $0.33$\,m, more than twice the
clearance of the route. Almost every flight the flow-based models lose ends inverted under the cascaded
geometric controller, and in the pilot of \autoref{sec:res:uav:controller} plans that invert under that
controller arrive under MuJoCo MPC, so this order may belong to the controller as much as to the models.}
\dataref{v3.87 (author: \enquote{also the S\_CURVE, one step straight, claim under simple data case it is not
so as we expected (though the conclusion not yet out, prefill part, mark as danger sign is ok)}).
\autoref{tab:uav-scurve} (R44a); tracking errors and the route's clearance:
\texttt{DA\_20260924\_scurve\_R44a\_raw\_grid.md}; lost flights of the flow-based models $57$, inverted $54$ (the
other three crossed after touching the floor). Demonstrations: \autoref{tab:uav-demos} and the waypoint-and-arc
construction of \autoref{sec:method:env:uav}. The expected order, the average-velocity objectives ahead at few
evaluations: \autoref{sec:bg:fewstep}.}

% ---- 6.4.1 prose and tab:summary-models ----
\subsection{Generative Models}\label{sec:res:summary:models}

On the two D3IL tasks the flow-based models Pareto-dominate the diffusion baseline, with the baseline's own
temporal U-Net (\autoref{tab:summary-models}). On D3IL-avoiding CI-MeanFM and FM at one network evaluation
hold its success with constraint satisfaction of $1.000$ in fewer control steps at a thirtieth of its time
per action. On D3IL-aligning analytic average-velocity matching under per-step projection keeps more contexts
free of violations than the baseline's best projected row, ends closer to the target and takes a third of its
time. The quadrotor scenes fly the same models on another plant. On UAV-pillars the operating points of
D3IL-avoiding keep, when flown, everything the declared constraints guarantee, and lose flights only to
collisions with pillars the constraints leave out. On UAV-corridor the flow-based models reach the end of the
corridor on every flight at a ninth to a twenty-fourth of the baseline's time and leave more violating steps,
a trade-off; on its straight-line demonstrations the three objectives cannot be told apart.

\begin{table}[htpb]
  \caption[Generative models across environments]{The conclusion of each environment on the generative models,
  as found in Sections~\ref{sec:res:state}--\ref{sec:res:uav}. Temporal U-Net of 3.96--4.0\,M parameters
  throughout; D3IL-aligning on training-set contexts. UAV-s-curve is pending.}
  \label{tab:summary-models}
  \centering
  \small
  \begin{tabularx}{\linewidth}{@{}p{2.2cm} L@{}}
    \toprule
    Environment & Conclusion \\
    \midrule
    D3IL-avoiding & Pareto dominance over the baseline: CI-MeanFM and FM at $\nfe=1$ hold S\&C $1.000$ in $59.2$
                    and $67.0$ control steps at $18.1$ and $17.3$\,ms, against the baseline's $1.000$ in $70.1$ at
                    $553.4$\,ms; MeanFM at $\nfe=1$ trades one episode in thirty ($0.967$) for the same saving \\
    D3IL-aligning & Pareto dominance over the baseline: MeanFM under per-step projection keeps nine contexts free of
                    violations against four and ends $0.179$ against $0.462$\,m from the target, at $266$ against
                    $809$\,ms; unprojected at $\nfe=20$ it closes $84\,\%$ of the distance, against $14$, $10$ and
                    $-2\,\%$ for CI-MeanFM, FM and the baseline \\
    UAV-pillars   & the D3IL-avoiding operating points flown by the quadrotor: the declared constraints transfer
                    exactly, with no violating step after projection at the same control steps and time per
                    action; S\&C falls from $0.967$--$1.000$ to $0.600$--$0.700$ only through collisions with
                    pillars outside the constraint sets, which end a quadrotor's flight \\
    UAV-corridor  & a trade-off with the baseline: every flight reaches the end of the corridor at a twenty-fourth
                    (tilt) and a ninth (hump) of its time per action, with more violating steps, at best $1.4$ and
                    $0.2$ against $0.8$ and $0.0$; on straight-line demonstrations the three objectives cannot be
                    told apart, and the budget decides \\
    \bottomrule
  \end{tabularx}
\end{table}



% ---- 6.4.3 conclusion paragraph ----
\subsection{Conclusion}\label{sec:res:summary:conclusion}

Read across the environments, the answer to the first research question is a Pareto dominance on the two
D3IL tasks: with the diffusion baseline's own temporal U-Net, the flow-based models hold or improve its
constraint satisfaction at a thirtieth (D3IL-avoiding) and a third (D3IL-aligning) of its time per action, in
fewer control steps on D3IL-avoiding and closer to the target on D3IL-aligning. On the quadrotor, UAV-pillars
flies the average-velocity operating points of D3IL-avoiding and shows that what the declared constraints
guarantee survives the change of plant, while collisions with obstacles the constraints leave out end a third
of the flights; on UAV-corridor the dominance becomes a trade-off, a small fraction of the baseline's time for
a few more violating steps. Which of the three objectives is used matters where the task separates them:
analytic average-velocity matching leads on D3IL-aligning, on D3IL-avoiding the budget and the selection rule
set their order, and on the straight lines of UAV-corridor they tie. The answer to the second research
question is set by the budget: endpoint projection beats the per-step projection of \ac{DPCC} at twenty
evaluations, on D3IL-aligning and over the hump of UAV-corridor, and has nothing to guide where one or two
evaluations solve the task. \autoref{tab:summary-combinations} collects the combination each environment
supports.


